\section{Discussion}

\subsection{What the results support}
On this benchmark a 4-billion-parameter model on a small GPU produced a usable plan for 18 of 20 questions when replanning was off, and the tools returned correct answers for all of them. No run timed out, the recovery suite passed, and every reported number is traceable to an evidence record. Almost all of the time is spent in the model, and the tools take seconds in total.

\subsection{What they do not support}
The best result, \NoRepQ{} of \NoRepN{} with replanning off, is not an unbiased estimate. We chose to report the replanning-off configuration as the better one after seeing both runs on the same questions, and the system had been changed after an earlier run on them. The default configuration, with replanning on, passes \MainQ{} of \MainN{}. The questions are also easy: one or two steps, with the tool nearly named by the wording. The result says little about longer chains or about questions where the first result should change what is asked next, which is where replanning is meant to help.

The baseline cannot be used as a competitor. It passes all twenty because several of its answers are assembled by question-specific code outside the tools. A fairer baseline would use the same answer extraction as the agent.

\subsection{Threats to validity}
\begin{itemize}
\item \textbf{Benchmark.} Twenty questions on one table, written by us. Several categories have two questions, so a category count is coarse. The question file and the expected answers were frozen before the runs, and we did not edit them.
\item \textbf{Post-freeze changes.} The tools, the step that extends plans, and the planner model were changed after an earlier run on these same questions, and the reported results are measured on the same questions. They are not a held-out measurement.
\item \textbf{Oracle.} The oracle imports no agent code, but it uses the same libraries as the tools. A defect in pandas that affected both would go unnoticed.
\item \textbf{Single run per question.} Decoding was deterministic (temperature 0, fixed seed), so a repeated run would give nearly the same outcome. Each question was run once, and we report no variance estimate.
\item \textbf{One hardware and model setup.} The main results use one model on one GPU. The comparison with three smaller models used replanning enabled only, and the 7B model does not fit in GPU memory.
\item \textbf{Answer rule.} The score reads the answer from the last executed task. This rule is what turned extra analyses into wrong answers in R1, so it interacts strongly with replanning. A different rule would give different numbers.
\item \textbf{Claim checking and narration.} The score never reads the report text. The claim checker and the narration step are tested by the recovery suite but are not measured on the 20 questions.
\end{itemize}

\subsection{Lessons}
Two checks were the most useful. Feeding hand-written plans through the scorer, without a model, showed which questions could be passed at all, and it separated failures of the interface from failures of planning. Comparing a run with replanning against one without it, question by question, showed that every difference came after the first plan and traced it to tasks added after a correct result. We also learned that a change made for one purpose can alter earlier results without any test noticing, which is a reason to run the whole question set again after each change.
